\section{Preserving the stronger degree-sensitive frontier}
\label{sec:degree-corollaries}

A comparison restricted to the generic $s^{9/10}$ degree interface would
miss useful stronger-premise bounds in the reviewed packets.  We record
two polynomial-space corollaries on the same normalized wire scale.  Both
retain the same permitted original entry-oracle cost as the generic
theorem.  Neither changes the unrestricted SAT logarithmic order.

\subsection{Root degree at most \texorpdfstring{$s^{3/5}$}{s to the power 3/5}}

Suppose the root degree is at most $\lceil s^{3/5}\rceil$.  Keep the
arities from Section~\ref{sec:variable-arity} and multiply every leading
retained-level coefficient $f_d$ by
\[
 \tau=4.
\]
Continue using $u_d=\lfloor\tau f_d\lambda\rfloor$ and
$a_d=a_{d-1}+u_d+\lceil6\log_2\lambda\rceil$.
Maintain the degree bound $\lceil s^{3/5}\rceil$ on the retained side.
The cumulative packet exponent is now
\[
 \tau A_\infty=\frac{499}{2500}=0.1996<\frac35,
\]
so every fresh packet tensor satisfies the maintained bound eventually.

The high-degree packing powers become $3/5+\tau h_d$, where
$h_d$ denotes the label-exponent part of the corresponding bound in
Section~\ref{sec:variable-arity}; their maximum is
\[
 \frac35+4\frac{2499}{25000}
 =\frac{6249}{6250}=0.99984<1.
\]
The root uses the doubled label exponent as before.  All pure label,
fresh-column and rounding powers are multiplied by $\tau$; their maximum
is at most
\[
 4\frac{499}{2000}
 =\frac{499}{500}=0.998<1.
\]
Thus all packing conditions retain fixed strict margins.  The same
coarse child shrinkage, retained-range ratio and depth argument apply.
Multiplying every $u_d$ by a fixed positive factor also leaves the
additive precision-gap proof unchanged.

The local costs are divided by $\sqrt\tau$.  Consequently this interface
has polynomial-space time bound
\[
 (s+\mathcal L)^{O(1)}
 2^{(C_{3/5}+o(1))s/\sqrt{\log_2s}},\qquad
 C_{3/5}=\frac{C_{\mathrm{multi}}}{2}
 =31.92870473555164\ldots.
\]
The integer coefficient $32$ is valid eventually: halve the earlier
rational upper bound on $C_{\mathrm{multi}}$ to obtain a bound below
$31.949$.  Packet-entry work is now at most
$\exp(s^{0.1996+o(1)})$, still lower order; all other resource arguments
remain as in the generic theorem.  An optional schedule with larger
power margins uses $\tau=39/10$ instead; it gives coefficient
$32.33545622568384\ldots$ and safe integer coefficient $33$.

This premise includes root degree $s^{1/2}\operatorname{polylog}(s)$ and
root degree $s^{0.51}$ for sufficiently large $s$.  Subject to their
common normalized precision, truncation and oracle hypotheses, it
therefore improves the respective $42.3726842\ldots$ and
$60.7069575\ldots$ degree-sensitive polynomial-space certificates from
the incoming high-moment and recursive packets.  This comparison is
between upper bounds on the same normalized network size.

\subsection{Root degree at most \texorpdfstring{$s^{1/12}$}{s to the power 1/12}}

For low-degree input it is better to compose the incoming asymmetric
packet's already proved low-degree schedule with the incoming full-block
Cauchy basis.  Maintain $\Delta=\lceil s^{1/12}\rceil$, and use
\[
 r=48,\qquad
 u_d=\left\lfloor\frac{\lambda}{14\,8^d}\right\rfloor,\qquad
 a_d=\sum_{j=0}^d\left(u_j+\lceil6\log_2\lambda\rceil\right).
\]
The leading cumulative packet exponent is at most $4/49$.  The exact
integer inequality $3^5<2^8$ yields $\log_2 48<28/5$.  Therefore the
symmetric degree-packing exponent is bounded by
\[
 2\frac4{49}\log_2 48+\frac1{12}
 <\frac{419}{420}<1,
\]
and the rounding exponent is bounded by
\[
 \left(1+2\log_2 48\right)\frac4{49}
 <\frac{244}{245}<1.
\]
The new packet degree has leading exponent $4/49<1/12$, so the same
small degree bound is preserved.  The coarse child ratio $1/12$ and
retained-range ratio $1/8$ give the usual noncircular depth bootstrap,
with a uniformly positive power of $\lambda$ in every active $u_d$.

The sharp child ratio is $1/24$.  The slope-four local coefficient is
\[
 K(48)=\frac{\pi}{\ln2}\sqrt{\frac{47}{24}}.
\]
Summing the convergent geometric series gives
\[
 C_{1/12}
 =\frac{K(48)\sqrt{14}}{1-\sqrt8/24}
 =\frac{\pi\sqrt{329/12}}
        {\ln2\,(1-\sqrt2/12)}
 =26.90232751306742\ldots.
\]
Thus the polynomial-space bound has eventual integer coefficient $27$.
The exact checker verifies this directly using rational enclosures.
This improves the prior polynomial-space coefficients
$31.0641321\ldots$ and $31.2364960\ldots$ on their overlapping low-degree
premises.  In particular, bounded-degree clause networks eventually
satisfy the $s^{1/12}$ premise.  Conversion to a SAT-variable coefficient
still requires accounting for the chosen encoding's wire budget.

These two corollaries are compositions with complete schedules and
resource accounting.  They do not imply numerical dominance between
networks measured in different wire parameters, and they do not compare
polynomial space with an unrestricted-memory claim as though the two
resource guarantees were identical.
